\documentclass[11pt]{article}
\usepackage[utf8]{inputenc}
\usepackage{amsmath,amssymb}
\usepackage{hyperref}
\title{Efficient Functional Mri Connectivity: A Resource-Aware Approach}
\author{Anonymous}
\date{}
\begin{document}
\maketitle

\begin{abstract}
This paper studies Functional Mri Connectivity. Grounded in a real, citable literature \cite{ref1,ref2,ref3,ref4,ref5,ref6,ref7,ref8},
we propose a focused method and report \textbf{illustrative} experiments.
All references below are real and verifiable (OpenAlex/arXiv); the reported
numbers are simulated to illustrate intended behaviour.
\end{abstract}

\section{Introduction}
Functional Mri Connectivity is an active research area. This work targets a concrete gap identified from
the recent literature and contributes a method together with an evaluation protocol.

\section{Related Work (real literature)}
The following works, retrieved from public bibliographic APIs, frame our study:
\begin{itemize}
\item Biswal et al. (1995) — "Functional connectivity in the motor cortex of resting human brain using echo‐planar mri", Magnetic Resonance in Medicine (cited 10134).
\item Yeo et al. (2011) — "The organization of the human cerebral cortex estimated by intrinsic functional connectivity", Journal of Neurophysiology (cited 9934).
\item Buckner et al. (2011) — "The organization of the human cerebellum estimated by intrinsic functional connectivity", Journal of Neurophysiology (cited 7910).
\item Power et al. (2011) — "Spurious but systematic correlations in functional connectivity MRI networks arise from subject motion", NeuroImage (cited 7910).
\item Seeley et al. (2007) — "Dissociable Intrinsic Connectivity Networks for Salience Processing and Executive Control", Journal of Neuroscience (cited 7583).
\item Power et al. (2011) — "Functional Network Organization of the Human Brain", Neuron (cited 4451).
\end{itemize}

\section{Method}
We reduce the dominant computational cost via adaptive resource allocation, activating only the components needed per input. We instantiate this idea for Functional Mri Connectivity and analyse its cost/quality trade-off.

\section{Experiments (Illustrative)}
\begin{center}
\begin{tabular}{lcc}
System & Primary Metric & Efficiency \\ \hline
Strong baseline & 0.812 & 1.00$\times$ \\
Ablation & 0.828 & 0.92$\times$ \\
\textbf{Ours} & \textbf{0.851} & \textbf{0.71$\times$}
\end{tabular}
\end{center}
\emph{Metrics simulated to illustrate the intended effect; not measured.}

\section{Conclusion}
We connected a real literature to a concrete method for Functional Mri Connectivity. Future work will
validate the illustrative results with real experiments.

\begin{thebibliography}{9}
\bibitem{ref1} Bharat B. Biswal, F. Zerrin Yetkin, Victor M. Haughton et al.. Functional connectivity in the motor cortex of resting human brain using echo‐planar mri. \emph{Magnetic Resonance in Medicine}, 1995. DOI:10.1002/mrm.1910340409
\bibitem{ref2} B.T. Thomas Yeo, Fenna M. Krienen, Jorge Sepulcre et al.. The organization of the human cerebral cortex estimated by intrinsic functional connectivity. \emph{Journal of Neurophysiology}, 2011. DOI:10.1152/jn.00338.2011
\bibitem{ref3} Randy L. Buckner, Fenna M. Krienen, Angela Castellanos et al.. The organization of the human cerebellum estimated by intrinsic functional connectivity. \emph{Journal of Neurophysiology}, 2011. DOI:10.1152/jn.00339.2011
\bibitem{ref4} Jonathan D. Power, Kelly A. Barnes, Abraham Z. Snyder et al.. Spurious but systematic correlations in functional connectivity MRI networks arise from subject motion. \emph{NeuroImage}, 2011. DOI:10.1016/j.neuroimage.2011.10.018
\bibitem{ref5} William W. Seeley, Vinod Menon, Alan F. Schatzberg et al.. Dissociable Intrinsic Connectivity Networks for Salience Processing and Executive Control. \emph{Journal of Neuroscience}, 2007. DOI:10.1523/jneurosci.5587-06.2007
\bibitem{ref6} Jonathan D. Power, Alexander L. Cohen, Steven M. Nelson et al.. Functional Network Organization of the Human Brain. \emph{Neuron}, 2011. DOI:10.1016/j.neuron.2011.09.006
\bibitem{ref7} Jessica S. Damoiseaux, Serge A.R.B. Rombouts, Frederik Barkhof et al.. Consistent resting-state networks across healthy subjects. \emph{Proceedings of the National Academy of Sciences}, 2006. DOI:10.1073/pnas.0601417103
\bibitem{ref8} Alexander Schaefer, Ru Kong, Evan M. Gordon et al.. Local-Global Parcellation of the Human Cerebral Cortex from Intrinsic Functional Connectivity MRI. \emph{Cerebral Cortex}, 2017. DOI:10.1093/cercor/bhx179
\end{thebibliography}
\end{document}
